\chapter{\MakeUppercase{Protocolo experimental}}
\label{cap:metodologia}


\section{Variáveis, tarefas avaliadas e unidade de análise}
\label{sec:tarefas-avaliadas}

A variável independente principal deste estudo corresponde à condição de
execução da elevação lateral bilateral. As variáveis dependentes incluem
os sinais inerciais produzidos pelos halteres, os resultados das etapas
de segmentação e classificação e a latência necessária para a geração
do resultado. A carga utilizada, o posicionamento inicial, o plano de
elevação, o número de repetições e os intervalos de descanso serão
controlados conforme o protocolo experimental.

A avaliação do HalterCheck foi organizada em tarefas consecutivas, mas
com diferentes papéis dentro da pesquisa. A aquisição e a sincronização
constituem requisitos técnicos para a obtenção dos sinais. A segmentação
identifica os intervalos correspondentes às repetições. A contagem é
obtida como resultado auxiliar da segmentação. A classificação dos
padrões de execução constitui a tarefa científica principal. A latência
é avaliada transversalmente para determinar a viabilidade de geração de
feedback em tempo quase real.

As unidades de análise variam conforme a tarefa. Na aquisição, a unidade
corresponde a uma amostra produzida por uma IMU e associada ao seu
dispositivo, sensor, número sequencial e instante de coleta. Na
segmentação, a unidade corresponde a um intervalo temporal contendo uma
repetição. Para a contagem, considera-se o conjunto de repetições de uma
série. Na classificação, a unidade principal é uma repetição bilateral
pareada, formada pelos sinais sincronizados dos halteres esquerdo e
direito. O Quadro~\ref{qua:tarefas-avaliadas} sintetiza as tarefas, seus papéis e
os respectivos critérios de avaliação.

\begin{quadro}[htb]
    \caption{Tarefas avaliadas no HalterCheck}
    \label{qua:tarefas-avaliadas}
    \centering
    \small
    \begin{tabular}{|p{0.19\textwidth}|p{0.25\textwidth}|p{0.45\textwidth}|}
        \hline
        \textbf{Tarefa} &
        \textbf{Papel no estudo} &
        \textbf{Critérios de avaliação} \\
        \hline

        Aquisição e sincronização &
        Requisito técnico &
        Frequência efetiva de amostragem, perdas, duplicações, pacotes
        fora de ordem, variação dos intervalos de amostragem e erro
        residual de sincronização entre os halteres.
        \\
        \hline

        Segmentação &
        Tarefa intermediária &
        Precisão, revocação e medida F1 na identificação das repetições,
        além do erro temporal na determinação de seus instantes inicial
        e final.
        \\
        \hline

        Contagem &
        Resultado auxiliar &
        Erro absoluto entre a quantidade prevista e a quantidade de
        repetições de referência, bem como a proporção de séries com
        contagem exata.
        \\
        \hline

        Classificação &
        Tarefa científica principal &
        Acurácia balanceada, precisão, revocação, medida F1 por classe,
        medida F1 macro e matrizes de confusão, utilizando validação com
        separação entre sessões.
        \\
        \hline

        Observabilidade postural &
        Análise associada à hipótese H5 &
        Desempenho relativo a uma referência externa de rotulagem e
        estabilidade da identificação dos erros posturais globais em
        sessões não utilizadas no treinamento.
        \\
        \hline

        Latência &
        Avaliação transversal &
        Latência mediana, percentil 95, valor máximo e proporção de
        repetições cujo resultado foi disponibilizado em até um segundo
        após o término do movimento.
        \\
        \hline
    \end{tabular}

    \begin{flushleft}
        \footnotesize Fonte: elaboração própria.
    \end{flushleft}
\end{quadro}

A tarefa de classificação será decomposta de acordo com as hipóteses da
pesquisa. Para H1, serão consideradas as condições de amplitude adequada,
posição inicial próxima às coxas, elevação insuficiente e elevação
excessiva. Para H2, serão avaliadas as condições de cadência de
referência, acelerada e reduzida. Para H3, serão comparadas execuções
simétricas e assimétricas. Para H4, serão avaliadas execuções com e sem
uso deliberado de impulso ou balanço.

Esses problemas serão avaliados separadamente, evitando que a expansão
do número de combinações entre erros produza classes com poucas
observações. As condições experimentais serão inicialmente executadas de
forma isolada, de modo que cada repetição possua um padrão principal
previamente definido.

A hipótese H5 será tratada como análise dos limites de observabilidade do
sistema, e não como uma classe principal de feedback. Os erros posturais
globais serão confirmados por uma referência externa utilizada
exclusivamente para rotulagem, sem que essa informação seja fornecida
como entrada aos métodos de classificação. Dessa forma, será possível
investigar até que ponto esses erros deixam manifestações indiretas nos
sinais dos halteres.

Os resultados relacionados à H5 serão interpretados como limitações da
configuração sensorial, do protocolo e dos métodos avaliados. Portanto,
um baixo desempenho não será apresentado como prova de impossibilidade
universal de identificação de erros posturais por halteres
instrumentados.

\subsection{Operacionalização das condições de amplitude}

Para avaliar a hipótese H1, foram utilizados como parâmetros o ângulo
inicial de elevação dos membros superiores, denotado por
$\alpha_{\mathrm{ini}}$, e o ângulo máximo atingido durante a repetição,
denotado por $\alpha_{\mathrm{max}}$. Esses ângulos serão determinados por
uma referência externa calibrada, utilizada exclusivamente para a
rotulagem e não como entrada dos métodos de classificação.

Foram estabelecidas as seguintes condições experimentais:

\begin{itemize}
    \item \textbf{amplitude adequada:}
    $10^\circ \leq \alpha_{\mathrm{ini}} \leq 15^\circ$ e
    $85^\circ \leq \alpha_{\mathrm{max}} \leq 95^\circ$;

    \item \textbf{posição inicial inadequada:}
    $\alpha_{\mathrm{ini}} \leq 5^\circ$ e
    $85^\circ \leq \alpha_{\mathrm{max}} \leq 95^\circ$;

    \item \textbf{elevação insuficiente:}
    $10^\circ \leq \alpha_{\mathrm{ini}} \leq 15^\circ$ e
    $\alpha_{\mathrm{max}} < 60^\circ$;

    \item \textbf{elevação excessiva:}
    $10^\circ \leq \alpha_{\mathrm{ini}} \leq 15^\circ$ e
    $\alpha_{\mathrm{max}} \geq 105^\circ$.
\end{itemize}

Os intervalos entre $5^\circ$ e $10^\circ$, entre $60^\circ$ e
$85^\circ$ e entre $95^\circ$ e $105^\circ$ constituem zonas de
transição. Repetições situadas nessas zonas serão registradas, mas não
serão utilizadas como exemplos induzidos das classes principais.

Durante a avaliação de H1, os dois membros deverão satisfazer a mesma
condição de amplitude. Repetições nas quais cada membro pertença a uma
classe diferente serão excluídas dessa análise, evitando a sobreposição
entre amplitude inadequada e assimetria.

As faixas estabelecidas constituem definições operacionais do protocolo
experimental e não pretendem estabelecer critérios clínicos universais
para a execução da elevação lateral.

\subsection{Operacionalização das condições de cadência}

Para avaliar a hipótese H2, a cadência foi definida pelas durações das
fases ascendente e descendente de cada repetição. Sejam
$T_{\mathrm{sub}}$ a duração da fase de subida,
$T_{\mathrm{desc}}$ a duração da fase de descida e
$T_{\mathrm{rep}} = T_{\mathrm{sub}} + T_{\mathrm{desc}}$ a duração
total da repetição.

Foram estabelecidas três condições:

\begin{itemize}
    \item \textbf{cadência de referência:}
    $1{,}75 \leq T_{\mathrm{sub}} \leq 2{,}25$ segundos,
    $1{,}75 \leq T_{\mathrm{desc}} \leq 2{,}25$ segundos e
    $3{,}5 \leq T_{\mathrm{rep}} \leq 4{,}5$ segundos;

    \item \textbf{cadência acelerada:}
    $0{,}75 \leq T_{\mathrm{sub}} \leq 1{,}25$ segundos,
    $0{,}75 \leq T_{\mathrm{desc}} \leq 1{,}25$ segundos e
    $1{,}5 \leq T_{\mathrm{rep}} \leq 2{,}5$ segundos;

    \item \textbf{cadência reduzida:}
    $2{,}75 \leq T_{\mathrm{sub}} \leq 3{,}25$ segundos,
    $2{,}75 \leq T_{\mathrm{desc}} \leq 3{,}25$ segundos e
    $5{,}5 \leq T_{\mathrm{rep}} \leq 6{,}5$ segundos.
\end{itemize}

As cadências serão orientadas por um metrônomo. Entretanto, a
classificação de referência de cada repetição será determinada pelas
durações efetivamente observadas, obtidas por uma referência temporal
externa, e não apenas pela condição que foi instruída.

Repetições cujas durações se encontrem entre as faixas estabelecidas
serão consideradas ambíguas e excluídas das classes induzidas. Também
serão excluídas repetições que apresentem pausa superior a
$0{,}25$ segundo no ponto de máxima elevação.

Nas condições de cadência, a amplitude deverá permanecer adequada e os
dois membros deverão executar a mesma condição temporal. Os valores
adotados representam ritmos experimentais de referência, não uma
prescrição universal de cadência correta.

\subsection{Operacionalização da assimetria entre os membros}

Para avaliar a hipótese H3, a assimetria foi delimitada à diferença
temporal entre os movimentos dos membros esquerdo e direito. Essa
delimitação evita sobreposição com os erros de amplitude avaliados em
H1.

Sejam $t_{\mathrm{ini},E}$ e $t_{\mathrm{ini},D}$ os instantes de início
do movimento esquerdo e direito, respectivamente, e
$t_{\mathrm{pico},E}$ e $t_{\mathrm{pico},D}$ os instantes em que cada
membro alcança a máxima elevação. Foram definidas as medidas:

\[
\Delta t_{\mathrm{ini}} =
\left|t_{\mathrm{ini},E}-t_{\mathrm{ini},D}\right|
\]

\[
\Delta t_{\mathrm{pico}} =
\left|t_{\mathrm{pico},E}-t_{\mathrm{pico},D}\right|.
\]

Uma repetição será considerada \textbf{temporalmente simétrica} quando:

\[
\Delta t_{\mathrm{ini}} \leq 0{,}15~\mathrm{s}
\quad \text{e} \quad
\Delta t_{\mathrm{pico}} \leq 0{,}15~\mathrm{s}.
\]

A condição de \textbf{assimetria temporal induzida} será executada com
um dos membros atrasado em relação ao outro. A repetição será incluída
nessa classe quando:

\[
\Delta t_{\mathrm{ini}} \geq 0{,}50~\mathrm{s}
\quad \text{e} \quad
\Delta t_{\mathrm{pico}} \geq 0{,}50~\mathrm{s}.
\]

Repetições com diferenças entre $0{,}15$ e $0{,}50$ segundo serão
consideradas ambíguas e não integrarão as classes induzidas. O lado
atrasado será registrado como metadado, e a quantidade de repetições
com atraso do membro esquerdo e do membro direito deverá ser
balanceada.

Durante essa condição, cada membro deverá manter individualmente a
amplitude adequada e a cadência de referência. Os instantes utilizados
para a rotulagem serão determinados por uma referência externa. A
comparação realizada pelo HalterCheck utilizará somente os sinais
sincronizados dos dois halteres.

Somente serão incluídas nessa análise sessões cujo erro residual de
sincronização entre os dispositivos seja inferior a um período de
amostragem.

\subsection{Padronização do plano de elevação}
\label{subsec:plano-elevacao}

Além da amplitude angular, foi padronizado o plano de execução do movimento. A elevação lateral foi realizada no plano escapular, com os membros superiores projetados aproximadamente $30^\circ$ à frente do
plano frontal do corpo. Essa orientação foi mantida em todas as condições experimentais, de modo que as alterações de amplitude, cadência e simetria fossem avaliadas sem a introdução deliberada de diferenças no
plano de elevação.

O plano de elevação não foi considerado uma variação de amplitude, mas uma característica distinta da trajetória do movimento. Sua possível identificação pelos sinais inerciais foi reservada para uma análise
exploratória, considerando as limitações do sistema sensorial utilizado para estimar orientação absoluta no plano horizontal.

\subsection{Critério de verificação das hipóteses H1, H2 e H3}

A capacidade de identificação das condições será avaliada em sessões
não utilizadas no treinamento dos modelos. As divisões dos dados serão
realizadas por sessão completa, evitando que repetições da mesma sessão
apareçam simultaneamente nos conjuntos de treinamento e teste.

Serão utilizadas como métricas principais a medida F1 macro e a
acurácia balanceada, acompanhadas das matrizes de confusão e das
medidas de precisão e revocação de cada classe. Regras manuais, Random
Forest, SVM e CNN serão avaliados utilizando as mesmas divisões dos
dados.

Será considerada evidência favorável à respectiva hipótese quando o
desempenho obtido em sessões não observadas durante o treinamento
superar o percentil 95 da distribuição de desempenho produzida pela
permutação aleatória dos rótulos. A estabilidade dos resultados entre
sessões também será apresentada.